\scanpage{64}
\mspage{31}
d'où, au moins au dessus de $\mathbf Z[\frac12]$
\[
(79)\quad\left\{\begin{aligned}
[M_{1,1}]&=\deleted{\unsure{\tfrac1{12}(T(1)-2)}}\frac12(T(1)-1)-\frac1{12}\\
[\widehat M_{1,1}]&=[M_{1,1}]+\underset{\substack{\Vert\\1/2}}{[\widehat M^\infty_{1,1}]}=\deleted{\unsure{\tfrac{T(1)}{12}+\tfrac13}}\frac12(T(1)-1)+\frac5{12}
\end{aligned}\right.
\]
Voici un autre calcul de $[M_{1,1}]$. On considère le morphisme modulaire
\[
M_{1,1}\longrightarrow\mathbf E^1
\]
\deleted{et} en découpant $\mathbf E^1$ suivant les deux sections $S_0$ (correspondant à la courbe elliptique harmonique de groupe d'automs $\mathbf Z/4\mathbf Z$) et $S_1$ \deleted{\unsure{$\simeq\mathbf Z$}} (courbe ell. \unsure{bonne} \unsure{avec un bon} gr. d'automs $\mathbf Z/6\mathbf Z$), qui ne se coupent qu'en les pts $2$ et $3$, qui correspondent (sauf erreur) à des groupes d'automs d'ordre $24$, et $12$. On a donc une partition
\[
\mathbf E^1=\deleted{\unsure{E}}U\amalg S'_0\amalg S'_1\cup a_2\cup a_3
\]
où \deleted{\unsure{$a_1,a_0$}}
\[
\left\{\begin{aligned}
U&=\mathbf E^1\smallsetminus(S_0\cup S_1)\\
\{\underset{\substack{|\\\text{car }2}}{a_2},\underset{\substack{|\\\text{car }3}}{a_3}\}&=S_0\cap S_1\\
S'_0&=S_0\smallsetminus\{a_2,a_3\},\quad S'_1=S_1\smallsetminus\{a_2,a_3\}
\end{aligned}\right.
\]
On trouve donc
\margin{63}
